% Portada de J2.5: puntos aleatorios en un cuadrado con un cuarto de círculo, y una campana de resultados.
\documentclass[tikz,border=0pt]{standalone}
\usepackage{resonancia}
\begin{document}
\begin{tikzpicture}
  \clip[rounded corners=6pt] (0,0) rectangle (12,3.6);
  \shade[top color=rAzul!16, bottom color=white] (0,0) rectangle (12,3.6);
  \draw[trazo, color=rGris, fill=white] (0.6,0.5) rectangle (3.6,3.2);
  \draw[color=rAzul, line width=1pt] (0.6,0.5) -- ++(2.7,0) arc (0:90:2.7) -- cycle;
  \foreach \x/\y in {0.9/0.8,1.4/2.6,2.2/1.1,2.9/0.7,3.3/1.8,1.1/1.5,2.6/2.8,1.8/1.9,3.1/3.0,0.8/2.9,2.0/0.6,1.5/0.9,3.4/2.4,2.4/2.1,1.2/3.0,2.8/1.5} {
    \pgfmathparse{(\x-0.6)^2+(\y-0.5)^2 <= 2.7^2 ? 1 : 0}
    \ifnum\pgfmathresult=1 \fill[rVerde] (\x,\y) circle (2pt); \else \fill[rNaranja] (\x,\y) circle (2pt); \fi
  }
  \draw[flecha] (3.9,1.9) -- (4.9,1.9);
  % campana
  \draw[eje] (5.3,0.8) -- (11.4,0.8);
  \draw[analisis, line width=1.4pt] plot[smooth, domain=5.5:11.2, samples=80] (\x, {0.8 + 2.2*exp(-((\x-8.35)^2)/(2*0.9))});
  \foreach \x/\h in {6.0/0.15, 6.5/0.3, 7.0/0.7, 7.5/1.2, 8.0/1.9, 8.5/2.1, 9.0/1.7, 9.5/1.1, 10.0/0.55, 10.5/0.25} {
    \fill[rAzul!35, draw=white] (\x-0.25,0.8) rectangle (\x+0.25,0.8+\h);
  }
  \node[etiqueta, color=rGris] at (8.35,0.45) {\small resultados de miles de simulaciones};
\end{tikzpicture}
\end{document}
